%% kapa: exam=1 round=27 year=2016 period=2 form=A subject=appraisal_law subject_ko=감정평가관계법규 seq=20 number=20
%% source: Q-Net 감정평가사 기출문제 (cst00309) artlSeq=5169402 file=extracted/1008/5169402/제27회 감정평가사 1차 문제/제27회 감정평가사 2교시문제(A형).pdf
%% conversion: kordoc 4.24.51 -> tools/gen1.py (auto); review=proofread-sonnet; flags=none
\begin{question}{20}
국유재산법령상 일반재산에 관한 설명으로 옳지 않은 것은?
\begin{choices}
  \item 일반재산은 대부 또는 처분할 수 있다.
  \item 총괄청은 일반재산의 관리ㆍ처분에 관한 사무의 일부를 위탁받을 수 있다.
  \item 일반재산인 토지의 대장가격이 3천만원 이상인 경우 처분예정가격은 하나의 감정평가법인의 평가액으로 한다.
  \item 일반재산은 개척사업을 시행하기 위하여 그 사업의 완성을 조건으로 대부ㆍ매각 또는 양여를 예약할 수 있다.
  \item 대부계약의 갱신을 받으려는 자는 대부기간이 끝나기 1개월 전에 중앙관서의 장등에 신청하여야 한다.
\end{choices}
\end{question}
\answer{3}
